\documentclass[11pt,a4paper]{article}
\usepackage[margin=24mm]{geometry}
\usepackage[T1]{fontenc}
\usepackage[utf8]{inputenc}
\usepackage{lmodern,amsmath,amssymb,amsthm,mathtools,booktabs,array,microtype}
\usepackage[dvipsnames]{xcolor}
\definecolor{ink}{HTML}{294B48}
\usepackage[colorlinks=true,linkcolor=ink,urlcolor=ink,citecolor=ink,pdftitle={Five-variable weak Markus--Yamabe: graded compression and affine-Hessian obstructions},pdfauthor={Anonymous}]{hyperref}
\usepackage{fancyhdr}
\pagestyle{fancy}\fancyhf{}\fancyhead[L]{\small Five-variable weak Markus--Yamabe}\fancyhead[R]{\small Research report}\fancyfoot[C]{\thepage}
\setlength{\headheight}{14pt}\setlength{\parskip}{5pt}\setlength{\parindent}{0pt}
\newtheorem{theorem}{Theorem}[section]\newtheorem{proposition}[theorem]{Proposition}\newtheorem{corollary}[theorem]{Corollary}
\theoremstyle{remark}\newtheorem{remark}[theorem]{Remark}
\newcommand{\R}{\mathbb R}\newcommand{\Q}{\mathbb Q}\newcommand{\Id}{\operatorname{Id}}\newcommand{\tr}{\operatorname{tr}}\newcommand{\Hess}{\operatorname{Hess}}
\title{\color{ink}Five-variable weak Markus--Yamabe:\\graded compression and affine-Hessian obstructions}
\author{Anonymous}\date{19 October 2026}
\begin{document}
\thispagestyle{plain}
\begin{center}
\vspace*{3mm}
{\LARGE\color{ink}Five-variable weak Markus--Yamabe:\\[2mm]
graded compression and affine-Hessian obstructions\par}
\vspace{6mm}
Anonymous\\[2mm]
19 October 2026\\[1mm]
\href{https://doi.org/10.5281/zenodo.23477226}{doi:10.5281/zenodo.23477226}
\end{center}
\vspace{2mm}
\begin{abstract}
We investigate the construction of a polynomial vector field $X:\R^5\to\R^5$
whose Jacobian has characteristic polynomial $(\lambda+1)^5$ everywhere and which
has two distinct real zeros. We do not obtain such a field. We prove that two
complete bigraded compression families built from a known six-dimensional
example contain no field with the requested characteristic polynomial. Exact
polynomial-identity certificates establish the exclusions. In the final chain
proof, eight affine-linear coefficient equations force the variables and a
ninth has exact nonzero residual. We also prove a uniform
Hessian-span theorem: for a specified linear form $A$ in three auxiliary
variables, every potential $\kappa A^2+D$, with $\kappa\ne0$ and arbitrary
$D$ affine in those variables, resists every invertible affine source--target
normalisation to constant Jacobian trace. A separate checker reconstructs the
Jacobians and replays the finite certificates. A tempting factor-sharing
compression retains three zeros but has an explicit positive Jacobian eigenvalue.
None of these conclusions rules out a five-dimensional counterexample outside
the stated architectures.
\end{abstract}

\section{Target, source boundary and status}
The target is an explicit real polynomial map satisfying
\begin{equation}\label{target}
 \det(\lambda I_5-DX(z))=(\lambda+1)^5\quad(z\in\R^5),
 \qquad X(z_1)=X(z_2)=0,\quad z_1\ne z_2.
\end{equation}
Writing $X=-\Id+H+\eta$, with constant $\eta$, makes the spectral requirement
\begin{equation}\label{nil}
 \det(LI_5-DH)=L^5.
\end{equation}
The constant $\eta$ affects fibers but not the Jacobian. Every exclusion below
concerns \eqref{nil}; it does not exclude a Hurwitz field with nonconstant
characteristic polynomial.

The source audit, updated on 19 October 2026, includes the
fourteen-dimensional chain construction of Casta\~neda, Honorato and
Valenzuela-Henr\'iquez \cite{chain}, the six-dimensional construction posted by
liqsweep \cite{six}, the five-variable Hessian potential of Meng and Yang
\cite{hessian}, a degree-bounded retargeting obstruction \cite{ksl}, and Liu's
Hesse-system formulation \cite{liu}. Earlier work already discusses failed chain
shortening and coefficient-span methods. The contributions here are the exact
displayed family exclusions, their transparent certificates, and the uniform
affine-trace obstruction below. This is a targeted source boundary, not an
exhaustive priority audit. The six-dimensional identities and its three listed
zeros were reproduced independently. The new results have not undergone external
specialist review or proof-assistant formalisation.

\textbf{Result status.} The requested five-dimensional construction remains
unresolved. The deliverables are exact obstruction theorems for specified
architectures, a rejected candidate, and reproducible exact evidence. Broader
uncertified searches are reported only as supplementary exploration.

\section{The baseline and the failed direct substitution}
For base variables $(x,y,z)$ put
\begin{align}
q&=y-6xz, & f&=x^2(3q+2xz),\\
g_2&=24xq(q+xz), &g_3&=2q(2q+3xz),\\
g_4&=12q^2 f, &g_6&=4q^3f.\nonumber
\end{align}
The remaining base polynomial is
\[
g_5=2xq^2(7q+6xz).
\]
The six-dimensional field in \cite{six} is
\begin{equation}\label{sixfield}
X_6=\begin{pmatrix}
-x+f\\1-y+a-g_2\\-z+b-g_3\\-a-g_4\\-b+c-g_5\\-c-g_6
\end{pmatrix}.
\end{equation}
Its characteristic polynomial is $(\lambda+1)^6$. The checker independently
expands this identity and verifies the zeros
\[
(0,1,-4,0,0,0),\quad(-1,-11,2,12,-6,4),\quad
(1,7,5/4,-3,-3/2,1/2).
\]
The proof mechanism in the source combines a constant determinant with a weighted
similarity identity. Its auxiliary variables have weights $2,2,3$ under the
base weights $(0,1,1)$.

The factorisation $g_6=(q/3)g_4$ suggests sharing an auxiliary variable. After
scaling the old $a$ by three, the tempting five-dimensional field is
\begin{equation}\label{naive}
X_{\rm try}=\begin{pmatrix}
-x+f\\1-y+3a-g_2\\-z+b-g_3\\-a-4q^2f\\-b+qa-g_5
\end{pmatrix}.
\end{equation}
It retains the three zeros
\[
(0,1,-4,0,0),\quad(-1,-11,2,4,-6),\quad
(1,7,5/4,-1,-3/2).
\]
It is nevertheless not a Hurwitz field.

\begin{proposition}\label{bad}
At $P=(1,0,0,-1,0)$ the field \eqref{naive} has a positive Jacobian eigenvalue.
More precisely,
\[
\det(\lambda I-DX_{\rm try}(P))
=(\lambda+1)^3\big((\lambda+1)^2-6\big).
\]
\end{proposition}
\begin{proof}
Direct differentiation gives
\[
DX_{\rm try}(P)=\begin{pmatrix}
-1&3&-16&0&0\\
0&-1&0&3&0\\
0&0&-1&0&1\\
0&0&0&-1&0\\
0&-1&6&0&-1
\end{pmatrix}.
\]
Expansion gives the stated polynomial. The eigenvalue $-1+\sqrt6$ is positive.
\end{proof}
This failure identifies an important distinction: eliminating auxiliary
variables in the equilibrium equations can preserve a fiber without preserving
the Jacobian away from that fiber. A determinant argument that evaluates an
auxiliary variable at its equilibrium value does not establish a global
Jacobian identity.

\section{Architectural coverage and relation to earlier obstructions}
The next two families are not arbitrary degree-bounded searches.  They are the
complete polynomial families with the displayed pairs of component gradings,
after fixing the inherited base terms and two nonzero auxiliary feedthroughs.
The feedthrough coefficients can be normalised to one by diagonal rescaling of
the auxiliary coordinates; this is a linear conjugacy and hence preserves the
Jacobian characteristic polynomial.  Cases in which either feedthrough is zero
are not covered by that normalisation.

\begin{proposition}\label{coverage}
For the parallel grading, the five components admit respectively
$2,5,4,8,8$ monomials.  Fixing the eight inherited base coefficients and two
feedthrough coefficients leaves $27-8-2=17$ free coefficients.  For the chain
grading, the corresponding counts are $2,4,4,7,11$; fixing the same data and
imposing the two trace relations leaves $28-8-2-2=16$ free coefficients.
Thus the formulae in Sections~\ref{parallelsection} and \ref{chainsection}
exhaust their stated bigraded architectures.
\end{proposition}
\begin{proof}
All variables except $x$ have positive first weight, while $x$ has second
degree one.  For each component, enumerate the finitely many exponents of the
positive-weight variables and solve the second-degree equation for the exponent
of $x$.  Retaining the nonnegative solutions gives the displayed lists and
counts.  The exact inventories are reproduced by
\texttt{discovery/grading\_audit.py}.
\end{proof}

Casta\~neda, Honorato and Valenzuela-Henr\'iquez already describe failed
shortening of their ordinary chain presentation \cite{chain}.  A later
degree-filtration result excludes a different retargeting strategy only for
target polynomials of degree at most eight \cite{ksl}.  Neither statement
contains the two bigraded families below.  Conversely, Proposition~\ref{coverage}
does not show that a successful five-dimensional construction must respect either
grading, retain the three fixed base components, or use these auxiliary
architectures.  For example, an $a^2$ term in the fourth parallel component has
first weight $4$ rather than the required component weight $3$ and is therefore
outside that family; so are changes to the fixed polynomial $f$ and branches
with a zero feedthrough.  The exclusions remove two natural complete search
spaces, not the global problem.

\section{Parallel auxiliaries: an exact exclusion}\label{parallelsection}
Consider the following family, with all coefficients arbitrary in a
characteristic-zero field:
\begin{equation}\label{parallel}
H=\begin{pmatrix}
f\\ a+\tau xb-g_2\\b-g_3\\-P+Aa+Bb\\-Q+Ca+Db
\end{pmatrix},
\end{equation}
where
\begin{align}
P&=\sum_{k=0}^3p_k x^{k+2}y^{3-k}z^k, &
Q&=\sum_{k=0}^3r_k x^{k+1}y^{3-k}z^k,\\
A&=\alpha xy+\beta x^2z, &
B&=\gamma x^2y+\delta x^3z,\\
C&=\epsilon y+\zeta xz, &
D&=\eta_1xy+\eta_2x^2z.
\end{align}
This is a 17-parameter family: eight cubic coefficients, eight auxiliary-coupling
coefficients and $\tau$. The first three nonlinear base terms are fixed to the
baseline. The family allows a cross-coupling $\tau xb$ in the second component.

Every displayed term has componentwise weighted degree $w_i+1$ for
$w=(0,1,1,2,2)$ and degree $\nu_i$ for
$\nu=(1,-1,-2,-1,-2)$. These two gradings explain the finite monomial lists;
they are explicit restrictions, not properties of all polynomial fields.

\begin{theorem}\label{paralleltheorem}
No choice of the 17 coefficients in \eqref{parallel} satisfies
$\det(LI-DH)=L^5$. Consequently no constant translation of $-\Id+H$ supplies
\eqref{target} in this family.
\end{theorem}
\begin{proof}
The fixed base terms satisfy
\[
f_x-(g_2)_y-(g_3)_z=0.
\]
Vanishing trace therefore forces $D=-A$, so
$\eta_1=-\alpha$ and $\eta_2=-\beta$. The coefficients linear in $a,b$ in the
next characteristic coefficient force
\begin{equation}\label{trace2relations}
\zeta=-\alpha-\tau\epsilon,\qquad
\gamma=\beta+\tau\alpha-3\tau.
\end{equation}
The checker reconstructs the unreduced family, checks its trace and the two
relevant coefficients of the second elementary symmetric function, and verifies
that these substitutions yield exactly the model used in the certificates.
Thus it suffices to examine the resulting 13-parameter family.

Differentiate \emph{before} setting $x=1$. Expand
\[
\det(LI-DH)|_{x=1}-L^5
\]
in $L,y,z,a,b$. There are 45 distinct nonzero coefficient equations in the
parameters. Thirteen do not involve $p_0,\ldots,p_3,r_0,\ldots,r_3$.
Denote these thirteen polynomials by $E_1,\ldots,E_{13}$.

The rational identity certificate
\texttt{cross\_pure\_certificate.json} proves
\begin{equation}\label{squarecert}
\beta^2,\ \epsilon^2\ \in (E_1,\ldots,E_{13}).
\end{equation}
It supplies 21 intermediate polynomial identities. The verifier reconstructs
the $E_i$ from the displayed Jacobian; it does not accept a separate list of
unverified equations as the input to the proof.

At a simultaneous zero of the coefficient equations, \eqref{squarecert} implies
$\beta=\epsilon=0$. One of the thirteen equations is
\[
-\alpha^2-\alpha\epsilon\tau-\beta\epsilon=0,
\]
so $\alpha=0$. A further coefficient, after these substitutions, is
\[
8\delta-128\tau=0.
\]
Every possible solution must consequently obey
\begin{equation}\label{forced}
\alpha=\beta=\epsilon=0,\qquad\delta=16\tau.
\end{equation}
The remaining auxiliary coupling is then $B=-\tau f$, and $A=C=D=0$.

Substitute \eqref{forced} into the 45 coefficient polynomials and discard zeros
and duplicates. The resulting 18 polynomials $F_1,\ldots,F_{18}$, in the nine
variables $p_0,\ldots,p_3,r_0,\ldots,r_3,\tau$, satisfy
\begin{equation}\label{unitcert}
1\in(F_1,\ldots,F_{18}).
\end{equation}
The certificate \texttt{cross\_final\_certificate.json} proves \eqref{unitcert}
through 31 further rational polynomial identities. A common zero would turn
\eqref{unitcert} into $1=0$, a contradiction.

Specialising $x=1$ loses no force in this exclusion: a global characteristic
identity would necessarily hold on that slice. No claim that checking one slice
establishes a global identity is used.
\end{proof}

\begin{remark}
The theorem excludes arbitrary values of the displayed coefficients, not merely
integer coefficients within a searched box. It does not exclude terms of other
bidegrees, changes to the fixed base terms, different dependence on the auxiliary
variables, or five-dimensional constructions unrelated to this baseline.
\end{remark}

\section{An unequal-weight chain: a second exclusion}\label{chainsection}
A different attempt shares the first auxiliary between the second and third
components, with auxiliary weights $2,3$. Let
\begin{align}
P&=\sum_{k=0}^3p_k x^{k+1}y^{3-k}z^k, &
Q&=\sum_{k=0}^4v_k x^{k+2}y^{4-k}z^k,\\
S&=\sum_{k=0}^2w_k x^{k+2}y^{2-k}z^k, &
R&=rxy+sx^2z.
\end{align}
After imposing the trace condition, the family is
\begin{equation}\label{chainfamily}
H=\begin{pmatrix}
f\\ \tau xa-g_2\\a-g_3\\b+P+Ra\\Q+Sa+jx^2a^2-Rb
\end{pmatrix}.
\end{equation}
Its sixteen coefficients are unrestricted. The weights are
$(0,1,1,2,3)$ and the second grading is $(1,-1,-2,-2,-2)$.
The quartic terms and the quadratic term $jx^2a^2$ are included; simply omitting
them would give a substantially narrower test.

\begin{theorem}\label{chaintheorem}
No member of \eqref{chainfamily} satisfies $\det(LI-DH)=L^5$.
\end{theorem}
\begin{proof}
Differentiate first and specialise $x=1$. Among the 38 distinct nonzero
characteristic coefficient polynomials are
\begin{align*}
E_0&=-2j-r\tau-s-3\tau,\\
E_1&=r\tau+s,\\
E_2&=-45r\tau+248r-8s\tau+42s,\\
E_3&=32j+16r\tau+3s\tau,\\
E_4&=-\tau(16r+3s).
\end{align*}
Their vanishing gives $32j=E_3+E_4=0$. Next
$E_0+E_1=-2j-3\tau$ gives $\tau=0$, then $E_1$ gives $s=0$, and finally
$E_2$ gives $r=0$.

After substituting $j=\tau=s=r=0$, 27 distinct nonzero coefficient equations
remain in the twelve coefficients of $P,Q,S$.  All are affine-linear.  Eight of
them successively give
\[
w_0=w_1=w_2=0,\qquad
p_1=132,\quad p_2=-744,\quad p_3=1392,\qquad
v_3=-1920,\quad v_4=3264.
\]
A further necessary equation is
\[
-93312w_2+25344v_4+34560v_3=0.
\]
The forced values make its left-hand side
\[
25344(3264)+34560(-1920)=16\,367\,616\ne0,
\]
which is the required contradiction.  The exact equation indices and a checked
rational linear combination are recorded by
\texttt{check\_chain\_short\_identity.py}.  The fuller thirteen-step
unit-ideal certificate is retained as redundant machine-checkable evidence.
\end{proof}

\section{A uniform affine-Hessian obstruction}
For variables $(x,y,a,b,c)$ set $u=1+xy$ and, for nonzero $r,t$, define
\[
A_{r,t}=a u^3+rxb u^2-tx^3c.
\]
Let $d_0,d_1,d_2,d_3\in\R[x,y]$ be arbitrary and put
\begin{equation}\label{generalphi}
\Phi=\kappa A_{r,t}^2+d_0(x,y)+a d_1(x,y)+b d_2(x,y)+c d_3(x,y),
\qquad \kappa rt\ne0.
\end{equation}
The coefficient-span object used below is the Hesse system in the terminology
of Liu \cite{liu}; that general linear-algebra criterion is not claimed as new.
The contribution is the following full-span calculation and its uniformity over
the perturbations in \eqref{generalphi}.

\begin{theorem}\label{hessiantheorem}
If a constant real $5\times5$ matrix $M$ satisfies
\[
\tr\bigl(M\Hess\Phi(w)\bigr)=\text{constant}\quad(w\in\R^5),
\]
then $M+M^{\mathsf T}=0$, and the constant is zero. Equivalently, the
nonconstant Hessian coefficient matrices of $\Phi$ span
$\operatorname{Sym}_5(\R)$.
\end{theorem}
\begin{proof}
Only the symmetric part $N=(M+M^{\mathsf T})/2$ contributes to the trace.
Split $N$ into the symmetric block on $(x,y)$, the mixed block between
$(x,y)$ and $(a,b,c)$, and the symmetric block on $(a,b,c)$. Compare
coefficients in descending auxiliary degree.

The coefficients of $b^2,a^2,x^2a^2$ first give
\[
2\kappa r^2 n_{xx}=0,\qquad
12\kappa n_{xy}=0,\qquad
30\kappa n_{yy}=0,
\]
so the base block vanishes.  Terms of auxiliary degree one then give
\[
4\kappa r n_{xa}=4\kappa r n_{xb}=24\kappa n_{ya}
=-12\kappa t n_{yc}=0
\]
and the two equations
\[
-12t n_{xc}+20r n_{yb}=0,\qquad
-16rt n_{xc}+16r^2 n_{yb}=0.
\]
Their coefficient determinant is $128r^2t\ne0$, so the mixed block vanishes.
Finally the coefficients of $xy,x,x^3,x^2,x^4,x^6$ give, respectively,
\[
12\kappa n_{aa},\quad4\kappa r n_{ab},\quad-4\kappa t n_{ac},\quad
2\kappa r^2 n_{bb},\quad-4\kappa rt n_{bc},\quad2\kappa t^2 n_{cc},
\]
and hence the last block vanishes.

The perturbation in \eqref{generalphi} cannot alter this triangular argument.
Its second base derivatives have auxiliary degree at most one, its mixed
derivatives have auxiliary degree zero, and its second auxiliary derivatives
vanish.  At each stage its possible contributions therefore lie in blocks
already eliminated or below the active diagonal block. Thus $N=0$.
Conversely, every skew-symmetric matrix pairs to zero with every symmetric
Hessian. The constant is therefore zero.

As an exact check, the selected $15\times15$ coefficient minor has determinant
\begin{equation}\label{hessdet}
-1\,304\,596\,316\,160\;\kappa^{15}r^{10}t^6,
\end{equation}
which is nonzero under the hypotheses.  At $(\kappa,r,t)=(1,3,1)$ its absolute
value is
\[
77\,035\,107\,872\,931\,840.
\]
\end{proof}

\begin{corollary}\label{affine}
For every $S,T\in\operatorname{GL}_5(\R)$ and $c_0,d_0'\in\R^5$, the map
\[
X(z)=T\nabla\Phi(Sz+c_0)+d_0'
\]
has nonconstant Jacobian trace. In particular it cannot satisfy \eqref{target}.
\end{corollary}
\begin{proof}
Cyclic invariance gives
\[
\tr DX(z)=\tr\bigl(ST\Hess\Phi(Sz+c_0)\bigr).
\]
Surjectivity of the source affine transformation and Theorem
\ref{hessiantheorem} would force the invertible matrix $ST$ to be
skew-symmetric. Every skew-symmetric matrix of odd order is singular, a
contradiction.
\end{proof}

The audited version of the Meng--Yang construction \cite{hessian} uses
$A=A_{3,1}$ and a potential whose remaining terms are affine in $(a,b,c)$;
every nonzero coefficient of $A^2$ is covered by \eqref{generalphi}.  This
formulation deliberately avoids relying on the numerical coefficient of the
source's linear-in-$A$ term.  More generally,
nonzero $r,t$ are coordinate rescalings of $r=3,t=1$, so this part records
coordinate robustness rather than a separate novelty claim.  The arbitrary
polynomials $d_i$ are the substantive uniform extension.  They need not preserve
a constant Hessian determinant or a gradient collision.  The theorem is only an
affine constant-trace obstruction; it does not exclude nonlinear transformations,
other leading quadratic forms or general five-dimensional Hurwitz maps.

\section{Certificate design and executed checks}
A certificate is a finite sequence of identities
\[
P_k=\sum_{i<k} c_{ki}\,m_{ki}P_i,
\]
where the $c_{ki}$ are exact rational numbers and the $m_{ki}$ are monomials.
The initial $P_i$ are coefficients reconstructed from the stated Jacobian.
The final targets are $\beta^2$, $\epsilon^2$ or $1$, as appropriate.
The checker only expands these identities. It need not implement, trust or
verify a Gr\"obner-basis search algorithm.

The discovery implementation uses SymPy's characteristic polynomial and a
certificate-producing Buchberger calculation. The independently written
standard-library checker uses sparse integer/rational polynomial arithmetic
and recursive determinant expansion. The checked bridges are therefore:
\begin{center}
explicit map $\longrightarrow$ differentiated matrix $\longrightarrow$
coefficient equations $\longrightarrow$ rational identities.
\end{center}
No sampled Jacobian evaluations are substituted for a polynomial identity.
The lone $x=1$ restriction is used only as a necessary condition in an
impossibility proof.

\begin{center}\small
\begin{tabular}{@{}p{.62\textwidth}p{.29\textwidth}@{}}\toprule
Obligation & Executed result\\\midrule
Six-dimensional characteristic identity & Zero residual polynomial\\
Three baseline zeros & Exact substitution passed\\
Naive compression & Positive-eigenvalue witness verified\\
Parallel compression family & Both staged certificates passed\\
Unequal-weight chain & Unit-ideal certificate passed\\
Short chain contradiction & Eight forcing equations plus a ninth residual equation; $16\,367\,616$\\
Original Hessian trace matrix & Submitted integer minor reproduced\\
Uniform Hessian theorem & Symbolic determinant \eqref{hessdet} verified\\
Grading completeness & Both monomial inventories reproduced\\
Four deliberate certificate corruptions & All rejected\\\bottomrule
\end{tabular}
\end{center}

The four mutation tests alter, respectively, a source equation, an intermediate
identity, the asserted final target, and the final derivation. All are rejected.
These tests detect several failure modes; they do not constitute a formal
verification of the checker itself.

To replay the proof-critical checks after extracting the bundle:
\begin{verbatim}
python3 -I verify.py --report /tmp/wmy5-verification.json
python3 -I test_mutations.py
python3 -I check_chain_short_identity.py
python3 discovery/grading_audit.py
python3 discovery/generalized_hessian_check.py
\end{verbatim}
The original verifier, mutation tests and short-chain check require only the
Python standard library. The two discovery audits require SymPy. Expected
terminal messages are
\texttt{ALL EXACT CHECKS PASSED} and \texttt{ALL CORRUPTIONS REJECTED}.
Optional SymPy discovery scripts can regenerate the rational certificates;
they are not needed to verify them. These scripts were also executed in a
separate copy of the bundle, and the standalone verifier accepted all three
regenerated certificates.

The report path above is outside the extracted package so that replay does not
overwrite the manifest-bound historical report.

\subsection*{Supplemental broader searches}
Exact Wolfram Language calculations returned the unit ideal for several larger
systems that allow changes to the lower base coefficients, as recorded in
\texttt{supplementary/CAS\_RESULTS.md}. The source script reconstructs those
systems. Small independent rational certificates for those broader systems are
\emph{not} included, and the standalone checker does not claim to verify them.
Some local Gr\"obner calculations exceeded an execution cap; those interruptions
supply no mathematical conclusion. These two statuses are kept separate.

\section{What remains open in this attempt}
The global construction problem \eqref{target} is not settled. The precise
conclusion is that a candidate cannot be obtained by selecting different
coefficients within the two displayed compression families. The entire potential
family \eqref{generalphi} also cannot be put into the requested constant-spectrum
form by invertible affine changes.

A construction outside these families could change the grading, break the
second homogeneity, use different nonlinear dependence on auxiliary variables,
or start from another Keller map. This report supplies no existence or
nonexistence claim for those possibilities. In particular it does not establish
that dimension six is minimal for polynomial constant-spectrum Hurwitz
counterexamples, or that all five-dimensional Hurwitz maps are injective.

\subsection*{Sanity checks}
The rejected candidate shows that preserving all three exhibited zeros is
insufficient: its Jacobian has an explicitly positive eigenvalue elsewhere.
For the Hessian obstruction the necessary target trace is $-5$, whereas every
constant trace obtainable by a constant left multiplier is zero. For the two
compression exclusions the final contradictions are exact ideal membership and
the explicit nonzero integer $16\,367\,616$, not numerical residuals or failures
to find coefficients.

\begin{thebibliography}{9}
\bibitem{six} blueberryvertigo [liqsweep]. (2026, September 6).
\emph{An integer polynomial Hurwitz counterexample in dimension six}
[Manuscript and verification code]. GitHub.
\url{https://github.com/blueberryvertigo/ArXiv-tryouts}.
README blob inspected: \texttt{002e811eeba67d7cade25d2d8f3003b9c5c38002}.

\bibitem{chain} Casta\~neda, \'A., Honorato, G., \& Valenzuela-Henr\'iquez, F.
(2026, August 5). \emph{The weak Markus--Yamabe conjecture fails in dimension 14}
[Preprint]. arXiv. \url{https://arxiv.org/abs/2608.05392}.

\bibitem{hessian} Meng, G., \& Yang, L. (2026, July 27).
\emph{A five-variable counterexample to the Hessian conjecture, and the
low-dimensional status of the Jacobian and Hessian conjectures}
[Preprint, version 2]. arXiv. \url{https://arxiv.org/abs/2607.22198v2}.

\bibitem{ksl} Korea Superintelligence Labs. (2026, August 30).
\emph{The degree filtration of the Alp\"oge subalgebra: Retargeting cannot
lower the Markus--Yamabe dimension below fourteen} [Preprint, version 1].
Ideosphere. \url{https://ideosphere.ai/papers/markus-yamabe-dim-v1.pdf}.

\bibitem{liu} Liu, H. (2026). \emph{On the Hessian conjecture in Lorentzian
signature: Constant pivots and Hesse systems} [Preprint]. arXiv.
\url{https://arxiv.org/abs/2608.19112}.
\end{thebibliography}
\end{document}
